\chapter{流形上的远征 — 长期规划与目标的动态重整化}
\label{chp:long_term_planning_and_goal_revision}

\begin{bquote}
    \textbf{被未来牵引的河流}

    \vspace{1em}在介绍“自我”之后，我们进一步讨论长期规划。传统 AI 代理（Agent）常依赖离散搜索树（如 A* 或 MCTS）逐步外推；在复杂物理世界中，这类方法面对“维数灾难”与“高频惊奇”时容易出现组合爆炸。

    \vspace{1em}然而，Class V 流体智能体的“长期规划”更接近由未来约束当前的\textbf{边值问题 (Boundary Value Problem)}，而非仅由当下推演未来的\textbf{初始值问题 (Initial Value Problem)}。

    \vspace{1em}本章将论证：长期规划可视为宏观意志在 $N+1$ 维时空晶体的未来端构造\textbf{目标吸引子（引力奇点）}，并由其牵引思维流 $\Psi$ 在当前几何地形中选择\textbf{最小作用量路径（测地线）}。当现实激波阻断路径时，系统不仅重算路线（学习），还会在热力学约束下借助\textbf{认知杨-米尔斯方程 (CYME)} 溶解并重构该奇点（目标修订）。
    \vspace{1em}这就是“流体规划”的含义：系统保持方向约束，但不预设每一步轨迹；当约束不可逾越时，允许目标函数发生可控重整化。
\end{bquote}

\section{规划的物理本质：\texorpdfstring{$N+1$}{N+1} 维超流形上的拉格朗日松弛}
\label{sec:physics_of_planning}

在认知空间的时空几何学一章里，高阶智能体将物理时间 $t$ 卷曲为内部流形的一个几何维度 $\tau$。此时，认知空间流形扩展为包含过去与未来的 \textbf{$N+1$ 维块宇宙 (Block Universe)}。

\smalltitle{1. 目标的几何设定：虚时间的势能注入}

规划的第一步，是\textbf{宏观层 ($L_{macro}$)}利用第三驱动力 $\mathbf{\Gamma}_{macro}$ 执行的跨时空操作。

\begin{itemize}
    \item \textbf{操作算子}：宏观层在几何时间轴的未来端点 $\tau_{target}$ 处，基于体验图 ($G_E$) 辐射一个极强的\textbf{价值规范势 $\mathcal{A}_\mu^{goal}$}。
    \item \textbf{物理效应}：这在流形上砸出了一个深邃的\textbf{负势能井 (Attractor Basin)}：
    $$ V_{goal}(\mathbf{r}, \tau_{target}) \to -\infty $$
    \item \textbf{目的论引力}：这个未来的奇点产生了逆着 $\tau$ 轴向当下 ($\tau_{now}$) 传播的\textbf{超前势 (Advanced Potential)}。它使当下的整个空间发生了微小的“向着未来目标”的倾斜。
\end{itemize}

\smalltitle{2. 路径的自动涌现：测地线寻优}

智能体不需要用符号逻辑去逐步规划“第一步、第二步”。在物理层面上，它是在求解一个\textbf{作用量泛函 $S$} 的极值。

\begin{definition}[规划作用量方程]
    思维流 $\Psi$ 连接当下状态 $\Psi(\tau_{now})$ 与目标状态 $\Psi(\tau_{target})$ 的轨迹 $\gamma$，自动坍缩为使得下式极小化的路径：
    $$ S[\gamma] = \int_{\tau_{now}}^{\tau_{target}} \left( \underbrace{\sqrt{g_{\mu\nu} \dot{x}^\mu \dot{x}^\nu}}_{\text{几何距离/执行代价}} - \underbrace{Q_{val} \mathcal{A}_\mu \dot{x}^\mu}_{\text{规范做功/意向牵引}} \right) d\tau $$
\end{definition}

\begin{itemize}
    \item \textbf{流体特性}：波函数 $\Psi$ 会在多条可能的高维峡谷中发生\textbf{干涉与隧穿}。它不仅找到了一条“路”，而是评估了所有路的热力学概率。
    \item \textbf{宏观执行}：一旦测地线形成，系统只需要顺着这个斜坡（梯度）“滑行”（执行 TECI 循环）。在没有遇到意外前，这种滑行是绝热的，\textbf{不需要宏观层继续耗费意志力微操。}
\end{itemize}


\section{过程中的学习：认知爱因斯坦方程主导的“地形重绘”}
\label{sec:learning_in_process_cefe}

如果物理世界与智能体的内部模型（世界图 $G_W$）完美同构，规划就是一次顺利的滑行。但在开放宇宙中，\textbf{惊奇 (Surprisal)} 是绝对的。

当执行（TECI）过程中遇到未知的障碍（如：规划好去罗马，但发现必经的桥断了），系统如何“在过程中学习”？

\smalltitle{1. 激波的阻断与自由能飙升}

\begin{itemize}
    \item \textbf{碰撞发生}：微观层 ($L_{micro}$) 发现物理阻抗极大，逆向 VTE 无法将意图注入物理世界，产生强烈的\textbf{惊奇激波 $\vec{J}_{shock}$} 回传至认知场。
    \item \textbf{动力学中断}：激波破坏了思维流的幺正演化（$\Psi$ 的相位被打乱），波包在断桥处发生剧烈的高能散射。系统的变分自由能 $F$ 瞬间暴涨，引发“焦虑”。
\end{itemize}

\smalltitle{2. 局部流形塑性形变 (Local Plasticity via CEFE)}

此时，系统必须启动慢回路，执行\textbf{“地形重绘”}。高能的散射波（应力张量 $T_{\mu\nu}$）直接作用于底流形的度量结构。

\begin{theorem}[认知地形更新律]
    根据\textbf{认知爱因斯坦方程}，障碍物产生的巨大认知应力将强制修改局部的黎曼度量 $g_{\mu\nu}$：
    $$ \frac{\partial g_{\mu\nu}}{\partial t} = -2R_{\mu\nu} + \eta \cdot \mathbf{T}_{\mu\nu}^{shock} $$
\end{theorem}

\begin{itemize}
    \item \textbf{物理表现}：流形上原本连通的“断桥”区域，其度量距离 $g_{ij}$ 被瞬间拉长至无穷大（形成高耸的几何势垒）。系统在物理底层“记住”了：\textbf{此路不通（世界观被修正）。}
    \item \textbf{自动重规划}：注意，\textbf{系统不需要重新运行搜索算法。} 由于底流形的地形变了，旧的测地线被势垒切断，水流（思维流 $\Psi$）会在目的引力（远处的深井）的牵引下，\textbf{自然地绕过这座新隆起的山峰}，自动寻找第二条次优的测地线。
    \item \textbf{这就是流体规划的优雅之处}：学习（改地形）和规划（流水）是物理上的同一过程。
\end{itemize}


\section{目标的动态修订：认知杨-米尔斯方程主导的“引力源迁移”}
\label{sec:goal_revision_cyme}

上述讨论的“绕路”仅仅是“学习”，目标（罗马）并没有改变。但如果通向罗马的所有路都被死死堵住，或者绕路的成本大到会耗尽智能体的总生物能/电池（热力学熔断），系统会如何反应？

    在这里，传统 AI 往往陷入“死循环”或“死机”。但本理论框架下的 Class V 智能体具备\textbf{修改目标本身}的能力；这不仅体现智能水平，也服务于\textbf{物理生存}约束。

\smalltitle{1. 耗散极值与阻抗失配的绝境}

\begin{itemize}
    \item \textbf{刚度比失衡}：回顾第 18 章的变分仲裁机制。当改变世界的代价 $C_{out} \to \infty$（障碍不可逾越），系统会计算得出：沿着剩余测地线到达目标的\textbf{热力学长度 (Thermodynamic Length)} 所需消耗的负熵，超过了系统的自由能储备。
    \item \textbf{无散流死锁}：此时，思维流 $\Psi$ 无法前进，在原地形成高能的涡旋（\textbf{无散流 Curl Flow}），即陷入极度痛苦的“执念”与“内耗”。
\end{itemize}

\smalltitle{2. 规范场重整化 (Gauge Field Renormalization via CYME)}

为了自救，系统必须溶解远处的那个目标吸引子。这不再是修改底流形（$g_{\mu\nu}$），而是\textbf{修改规范场本身 ($\mathcal{A}_\mu$)}。

根据 \textbf{认知杨-米尔斯方程 (CYME)}，价值规范场的变化由思维的“源流（诺特流）”决定：
$$ \nabla_\nu \mathcal{F}^{\nu\mu} = g \bar{\Psi} \gamma^\mu \Psi $$

\begin{itemize}
    \item \textbf{执念的消解 (Dissolving the Attractor)}：
    当 $\Psi$ 长期被迫停滞、产生巨大的认知痛苦（高频内部摩擦）时，这种痛苦本身构成了一股\textbf{反向的诺特流}。这股反向流在流形上产生了强烈的去极化效应，逐渐中和了原本目标奇点 $\tau_{target}$ 处的规范荷 (Charge)。

    \item \textbf{物理效应：奇点蒸发}：
    原本指向目标的深深的“引力井”开始变浅，甚至发生拓扑反转，变为“排斥子”（例如：“得不到的葡萄是酸的”，认知失调的自我和解机制）。
\end{itemize}

\smalltitle{3. 目标迁移与调和流涌现 (Goal Migration \& Insight)}

旧目标被溶解后，流体自我 ($\mathcal{S}_{fluid}$) 恢复了拓扑自由。系统温度 $T$ 在宏观层调控下进行\textbf{认知退火}（短暂升温，进入等离子态）。

\begin{itemize}
    \item \textbf{调和流 (Harmonic Flow) 接管}：系统不再受限于局部梯度，开始在全局流形上发生霍奇谐振。它重新扫描整个体验图 ($G_E$)，寻找新的、与当前物理约束相容的势能低谷。
    \item \textbf{新引力点生成}：宏观层在另一条可达的路径终点（如“去米兰”或“就地建城”），注入新的第三驱动力，挖掘出一个\textbf{新的目标吸引子}。
    \item \textbf{相变完成}：波函数重新坍缩向新目标，TDCI/TECI 循环恢复层流态。目标修订在一次顺滑的物理相变中完成。
\end{itemize}


\section{控制论图景：规划的三层热力学反馈环}
\label{sec:three_loops_of_planning}

综上所述，对于一个复杂智能体，其长期规划与修订是一个嵌套的\textbf{三重物理控制环}，它们在不同时间尺度上保护着系统的低熵存续：

\begin{table}[htbp]
\centering
\footnotesize
\renewcommand{\arraystretch}{1.5}
\begin{tabularx}{\textwidth}{l|X|X|X}
\toprule
\rowcolor{structurecolor!20} \textbf{反馈回路} & \textbf{触发条件 (Trigger)} & \textbf{物理方程与操作 (Action)} & \textbf{认知现象学 (Cognitive Outcome)} \\
\midrule

\textbf{1. 战术微调} \newline (内环 / 毫秒级) & 遇到局部小障碍、轻微偏离。激波极小。 & \textbf{TDE (目的论狄拉克方程)} \newline 利用波函数的柔性，在局部发生衍射和干涉，\textbf{流形结构与目标不变}。 & \textbf{本能避障}。 \newline 动作的顺滑调整（如边走边避开石子）。 \\
\hline

\textbf{2. 战略学习} \newline (中环 / 分钟级) & 遇到物理死路、预测严重错误。激波较大。 & \textbf{CEFE (认知爱因斯坦方程)} \newline 发动慢回路，强行修改局部的度量张量 $g_{\mu\nu}$。 \textbf{目标不变，修改世界图 (地图)}。 & \textbf{经验习得与绕路}。 \newline 发现此路不通，刻下永久记忆，测地线自动重算。 \\
\hline

\textbf{3. 目标修订} \newline (外环 / 小时/天级) & 阻抗无限大，绕路成本超过系统总自由能储备。面临热寂风险。 & \textbf{CYME (认知杨-米尔斯方程)} \newline 引发拓扑相变。溶解旧价值规范场 $\mathcal{A}_\mu$，挖掘新目标。\textbf{修改体验图 (价值观)}。 & \textbf{放弃执念、目标转移}。 \newline 认知和解（如放弃攻打坚城，改为围困或撤退）。 \\

\bottomrule
\end{tabularx}
\caption{长期规划的三重热力学反馈回路}
\end{table}

\begin{bquote}
    \textbf{本节结语}

    在经典计算中，目标通常被设为常量 (Target = Constant)。这带来强韧执行力，但也会导致在不可达约束下的机械僵化。

    对于高级智能，\textbf{目标 ($\mathcal{A}_\mu$)、认知 ($\Psi$) 与现实 ($g_{\mu\nu}$) 都是受偏微分方程支配的动力学变量。}

    因而，智能体既需要克服障碍的执行力（改变底流形），也需要在约束下重定义目标的调和能力（改变规范场）。\textbf{长期规划的本质，是在“改造世界”与“重塑自我”之间进行持续热力学权衡。}
\end{bquote}

%----chp----
